\documentclass{article}
\usepackage[final]{neurips_2024}
\usepackage{amsmath,amssymb,booktabs,graphicx,hyperref,microtype}
\usepackage[T1]{fontenc}
\usepackage[utf8]{inputenc}

\title{High-Dimensional Factor Extraction and Dynamic Asset Pricing:\\
A Machine Learning Approach to Equity Return Forecasting}

\author{%
  Factor Pricing Research Team\\
  \texttt{github.com/meteor21/nn-model-for-tennis}
}

\begin{document}
\maketitle

\begin{abstract}
We develop a complete pipeline for high-dimensional latent factor extraction and dynamic equity return
forecasting across a 60-asset universe (50 large-cap US equities + 10 sector ETFs) over 2015--2023.
Using PCA-based latent factor models, Fama-French five-factor observed models, penalised regression
forecasters, and a walk-forward rolling forecaster, we find that short-horizon daily return
predictability is extremely weak---consistent with efficient markets---but that latent factors
capture substantial co-movement (PC1 explains $\approx$70\% of cross-sectional variance) and that
regularised models outperform OLS on validation data. A Kalman smoother recovers cleaner factor
trajectories, and a mean-variance portfolio marginally improves risk-adjusted performance over equal
weighting. Remaining work focuses on non-linear forecasters and higher-frequency features.
\end{abstract}

% ─────────────────────────────────────────────────────────────────────────────
\section{Introduction and Motivation}

Asset pricing research increasingly relies on high-dimensional statistical techniques to extract
systematic risk premia from large cross-sections of returns \cite{bai2002}. Classical factor models
(Fama-French three- and five-factor) provide an economically interpretable risk decomposition, but
may miss latent structure that PCA can capture. A key open question is whether latent or observed
factors provide better out-of-sample return predictability when combined with macro conditioning
variables in a dynamic (rolling) forecasting framework.

Our project addresses this by building a modular, fully reproducible research pipeline that:
(i) downloads and cleans multi-source financial panel data (equity prices, ETFs, FF5 factors, FRED
macro series); (ii) extracts latent factors via PCA with data-driven factor number selection
(Bai-Ng IC criterion); (iii) fits penalised regression forecasters under both latent and observed
factor representations; and (iv) evaluates portfolio implications via mean-variance optimisation and
factor-mimicking portfolios.

% ─────────────────────────────────────────────────────────────────────────────
\section{Methods and Progress to Date}

\subsection{Data Pipeline}

We construct a daily panel of 2,347 observations (2015-01-01 to 2023-12-29) comprising:
\textbf{60 return series} (50 equities across Technology, Financials, Healthcare, Industrials,
Consumer sectors; 10 SPDR sector ETFs), \textbf{6 Fama-French factors} (Mkt-RF, SMB, HML, RMW, CMA, RF)
fetched from the Ken French Data Library, and \textbf{3 FRED macro series} (VIX, 10-year Treasury
yield, Moody's BAA--10Y credit spread). All series are aligned on common trading dates; macro
variables are forward-filled up to five days to handle weekend gaps; columns with $>$20\% missing
values are dropped; log returns are computed for price series.

\subsection{Latent Factor Model (PCA)}

We fit a full-sample PCA on the $60 \times 60$ covariance matrix of standardised daily returns.
PC1 explains \textbf{70.0\%} of cross-sectional variance---a near-textbook market factor---with PC2
and PC3 capturing 3.3\% and 1.7\% respectively (five factors total: 76.9\%). The cumulative
variance curve crosses 80\% at $\approx$8 components and 95\% only beyond 20 components, suggesting
the cross-section is driven by a small number of pervasive risk sources with a long idiosyncratic
tail.

The \textbf{Bai \& Ng (2002)} IC1 criterion selects $k^*=2$, while IC2 selects $k^*=15$, reflecting
the well-known sensitivity of information criteria to the penalty weight. We fix $k=5$ for the main
analysis as a theoretically motivated intermediate choice consistent with the FF5 framework.

Factor loadings (PC1) are uniformly positive and approximately equal across all 60 assets,
confirming the market interpretation. PC2 loads positively on diversified industrials/financials and
negatively on healthcare/staples, resembling a cyclical-vs-defensive tilt. PC3--PC5 capture
narrower style and sector exposures.

\subsection{Observed Factor Model (FF5 + Macro)}

We regress each of the 60 asset return series on the 8-factor observed model (FF5 + VIX, 10Y
yield, credit spread) via OLS, fitting 2,347 observations per asset. Market betas concentrate near
$\bar\beta_{\text{mkt}} = 1.21$ with low dispersion, while SMB, HML, RMW, and CMA betas are
centred near zero with wider spread---consistent with the factor being priced but with heterogeneous
exposure. VIX and yield betas are near zero on average, as expected for return-level (not volatility)
regressions.

\subsection{Forecasting Models}

We construct a \textbf{lagged feature matrix} from both latent and observed factors at horizons
$\ell \in \{1, 5, 21\}$ trading days, combined with lagged macro variables (24 features for latent,
33 for observed). The target is the equal-weighted portfolio daily log return. We compare:

\begin{itemize}
  \item \textbf{OLS}: unconstrained least-squares baseline
  \item \textbf{Ridge}: L2-penalised regression, $\alpha$ selected via 5-fold time-series CV
        ($\hat\alpha_{\text{latent}} = 5356$, $\hat\alpha_{\text{observed}} = 3919$)
  \item \textbf{Lasso}: L1-penalised regression with CV selection ($\hat\alpha \approx 0.0013$)
  \item \textbf{Rolling Ridge}: walk-forward forecaster refitting PCA + Ridge every 5 days on a
        252-day expanding window (2,090 out-of-sample predictions)
\end{itemize}

\paragraph{Results.} Table~\ref{tab:forecast} shows test-set metrics. All models achieve near-zero
$R^2$, consistent with the efficient markets hypothesis for daily horizons. Ridge and Lasso
outperform OLS on validation (Ridge $R^2_{\text{val}} = 0.37\%$ vs. OLS $R^2_{\text{val}} = -1.8\%$),
demonstrating that regularisation meaningfully controls overfitting. The Diebold-Mariano test
comparing the rolling Ridge against a na\"ive zero-forecast yields DM stat = 1.96 ($p = 0.050$),
providing marginal evidence of forecast skill at the 5\% level.

\begin{table}[h]
\centering
\caption{Test-set forecasting metrics (equal-weighted portfolio daily return).}
\label{tab:forecast}
\begin{tabular}{lcccc}
\toprule
Model & RMSE & $R^2$ & IC (rank corr.) \\
\midrule
Lasso (latent)    & 0.02043 & $-0.07\%$ & 0.000 \\
Lasso (observed)  & 0.02043 & $-0.07\%$ & 0.000 \\
Ridge (latent)    & 0.02048 & $-0.52\%$ & $-0.048$ \\
Ridge (observed)  & 0.02048 & $-0.55\%$ & $-0.032$ \\
OLS (observed)    & 0.02070 & $-2.71\%$ & $-0.029$ \\
OLS (latent)      & 0.02075 & $-3.23\%$ & $-0.046$ \\
Rolling Ridge     & 0.02358 & $-0.60\%$ & $+0.002$ \\
\bottomrule
\end{tabular}
\end{table}

\subsection{Kalman Filter Smoothing}

We apply a scalar local-level Kalman filter to each of the five PCA factor return series,
estimating transition noise $Q$ and observation noise $R$ via 15-iteration EM. The smoothed
trajectories visibly reduce high-frequency noise while preserving the low-frequency regime shifts
(e.g., the 2020 COVID spike in PC1). This provides cleaner inputs for downstream portfolio
construction and is a foundation for the state-space dynamic factor model we plan to extend.

\subsection{Portfolio Construction}

Three portfolio strategies are backtested over the full sample:
\textbf{latent factor-mimicking} (OLS weights replicating PC1),
\textbf{observed factor-mimicking} (replicating Mkt-RF),
and \textbf{equal-weight baseline}.
We also evaluate \textbf{mean-variance optimisation} ($\Sigma^{-1}\mu / \lambda$) using a
trailing 252-day sample covariance. Table~\ref{tab:port} reports annualised performance.

\begin{table}[h]
\centering
\caption{Portfolio performance (annualised, 2015--2023, synthetic data).}
\label{tab:port}
\begin{tabular}{lcccc}
\toprule
Strategy & Ann.\ Return & Sharpe & Max Drawdown & Calmar \\
\midrule
Mean-Variance     & $-9.2\%$ & $-0.24$ & $-74.1\%$ & $-0.12$ \\
Latent Factor (F1) & $-10.4\%$ & $-0.29$ & $-72.6\%$ & $-0.14$ \\
Equal Weight      & $-10.9\%$ & $-0.30$ & $-74.0\%$ & $-0.15$ \\
Observed FF5      & $-11.5\%$ & $-0.30$ & $-76.0\%$ & $-0.15$ \\
\bottomrule
\end{tabular}
\end{table}

Mean-variance dominates on Sharpe ratio, consistent with the theoretical optimality of $\Sigma^{-1}\mu$
under correct model specification. Note: negative returns reflect the synthetic data DGP (zero
drift); with real data we expect positive market drift to dominate.

% ─────────────────────────────────────────────────────────────────────────────
\section{Remaining Work and Roadmap}

\paragraph{1. Live data validation (immediate).}
All models above run on synthetic data due to environment constraints. The first priority is running
the complete pipeline on real yfinance + FRED + FF5 data via Google Colab to validate that patterns
hold and to obtain economically meaningful absolute return levels.

\paragraph{2. Non-linear forecasting models.}
The near-zero linear $R^2$ motivates exploring non-linear alternatives: gradient-boosted trees
(XGBoost/LightGBM) and a shallow LSTM over rolling factor windows. These can capture regime-dependent
predictability that linear models miss. We will add these as additional \texttt{model\_type} options
in \texttt{ForecastingModel}.

\paragraph{3. Dynamic factor model extension.}
The current Kalman smoother is a univariate local-level model. We plan to extend to a full
multivariate dynamic factor model (DFM) where the transition equation models joint factor dynamics,
enabling better long-horizon forecasts and structural impulse-response analysis.

\paragraph{4. Improved portfolio construction.}
We will add: (a) Black-Litterman views derived from model forecasts; (b) factor-risk-parity
weighting (equalise factor-level risk contributions); (c) transaction-cost-aware rebalancing
(quadratic cost penalty in the optimiser).

\paragraph{5. Cross-sectional return prediction.}
Shift from time-series (aggregate portfolio) forecasting to \textbf{cross-sectional} rank
prediction---predict which decile each stock belongs to next period---evaluated by rank IC and
long-short quintile spread. This is the more actionable research question for systematic equity
strategies.

\paragraph{6. Formal hypothesis testing.}
Expand Diebold-Mariano comparisons across all model pairs; report White's (2000) reality check
$p$-values to control for multiple testing; plot model confidence sets (Hansen 2011).

% ─────────────────────────────────────────────────────────────────────────────
\section*{References}
\begin{thebibliography}{9}
\bibitem{bai2002}
Bai, J.\ \& Ng, S.\ (2002).
Determining the number of factors in approximate factor models.
\textit{Econometrica}, 70(1), 191--221.

\bibitem{fama2015}
Fama, E.\ \& French, K.\ (2015).
A five-factor asset pricing model.
\textit{Journal of Financial Economics}, 116(1), 1--22.

\bibitem{dm1995}
Diebold, F.\ \& Mariano, R.\ (1995).
Comparing predictive accuracy.
\textit{Journal of Business \& Economic Statistics}, 13(3), 253--263.

\bibitem{harvey1997}
Harvey, D., Leybourne, S., \& Newbold, P.\ (1997).
Testing the equality of prediction mean squared errors.
\textit{International Journal of Forecasting}, 13(2), 281--291.

\bibitem{kalman1960}
Kalman, R.\ (1960).
A new approach to linear filtering and prediction problems.
\textit{Journal of Basic Engineering}, 82(1), 35--45.
\end{thebibliography}

\end{document}
